\begin{problem}[Local rings are invariant under shrinking]
Let $f\notin\mathfrak p$.  Prove
\[
(A_f)_{\mathfrak pA_f}\cong A_{\mathfrak p}
\]
and show the residue fields agree.  Explain why this is the algebraic form of saying that a germ depends only on arbitrarily small neighborhoods.
\end{problem}

\paragraph{Why this problem matters.}
This theorem is the bridge from affine localization to the later sheaf-theoretic notion of a stalk.

\paragraph{Useful ideas before starting.}
In $A_{\mathfrak p}$ the element $f$ is already invertible.

\paragraph{Guided hints.}
Use the universal property twice.

\begin{solution}
Since $f\notin\mathfrak p$, the map $A\to A_{\mathfrak p}$ factors through $A_f$.  Localizing $A_f$ at $\mathfrak pA_f$ then inverts all images of elements of $A\setminus\mathfrak p$, so it factors through $A_{\mathfrak p}$.  Conversely the canonical maps from $A$ give a map $A_{\mathfrak p}\to(A_f)_{\mathfrak pA_f}$.  The two are inverse by uniqueness of localization maps.  Quotienting by maximal ideals gives equal residue fields.
\end{solution}

\paragraph{What happened mathematically.}
Local data at a point is unaffected by discarding parts of the space that do not contain the point.

\paragraph{Extensions.}
Replace a single principal open by any affine neighborhood containing the point once VI/17 is available.

\paragraph{Provenance.}
New canonical synthesis problem based on the local-ring and residue-field corpus.
